% % arara: xelatex: {
% % arara: --> interaction: nonstopmode,
% % arara: --> options: ['-file-line-error','-max-print-line=200'],
% % arara: --> shell: yes,
% % arara: --> synctex: yes,
% % arara: --> }
% arara: lualatex
% arara: makeglossaries
%% arara: makeindex

\documentclass[letterpaper,10pt,twoside,final,es-MX]{article}
%%%%%%%%%%%%%%%%%%%%%%%%%%% Control de versiónes %%%%%%%%%%%%%%%%%%%%%%%%%%%%%%%%%%
\usepackage{changelog}

%%%%%%%%%%%%%%%%%%%%%%%%%%% ESPECIALES %%%%%%%%%%%%%%%%%%%%%%%%%%%%%%%%%%
\special{dvipdfmx:config z 0}	% Desactivar compresión
\special{pdf:minorversion 7}

%%%%%%%%%%%%%%%%%%%%%%%%%%% PREAMBULO %%%%%%%%%%%%%%%%%%%%%%%%%%%%%%%%%%
\usepackage{fontspec}
	\setmainfont{Carlito}

\usepackage{polyglossia}
	\setdefaultlanguage[variant=mexican,spanishoperators=all]{spanish}
	%\AtBeginDocument{\renewcommand\babelname{spanish}}
\usepackage{tracklang}
\SetCurrentTrackedDialect{es-MX}
\CurrentTrackedDialectVariant{}

\usepackage[toc=true,postpunc=comma,showtargets=annoteleft,style=tree,xindy]{glossaries-extra}
\makeglossaries
\setabbreviationstyle[acronym]{long-short}
\loadglsentries[main]{../glosario-def.tex}

\usepackage{imakeidx}
\makeindex[columns=2]

\usepackage[multiple]{footmisc}
\usepackage{microtype}
\usepackage{tabularray}
\usepackage{graphicx}
\usepackage{hyperref}
\usepackage{cleveref}
\usepackage{array}
\usepackage{lastpage}
\usepackage{siunitx}
	\sisetup{detect-all}
\usepackage{color}
\usepackage{tabularray}
\UseTblrLibrary{booktabs,siunitx}
\usepackage{xhfill}
%\usepackage{pdfpages}
%\usepackage{lipsum}

\usepackage[margin=2cm,bmargin=2cm,tmargin=5cm,headheight=97pt]{geometry}
%\addtolength{\topmargin}{-12pt}
	%\usepackage{geometry}
	%\setlength{\parindent}{0.95cm}
	%\renewcommand{\labelitemi}{\textemdash}

\usepackage{enumitem}
	\setlist[1]{labelindent=\parindent} % < Usually a good idea
	\setlist[itemize]{leftmargin=*}
	\setlist[itemize,1]{label=---}
	\setlist{noitemsep}

\usepackage{fancyhdr}
	\renewcommand{\headrulewidth}{0pt}
	\fancyhead[CE,CO,LE,LO,RE,RO]{}

	\newcommand{\Codigo}{}
	\newcommand{\FechaPub}{}
	\newcommand{\Titulo}{}
	\newcommand{\MayorVer}{}
	\newcommand{\MenorVer}{}
	\newcommand{\Edit}{\MayorVer.\MenorVer}

	\newcommand{\email}[1]{\href{mailto:#1}{#1}}

\fancypagestyle{formato}{%
    \fancyhead[C]{%
	    \begin{tblr}{%	
            colspec = {|X[c,m]|X[c,m]|X[c,m]|X[c,m]|X[c,m]|X[c,m]|},
            }
            \hline
            \SetCell[c=2]{c} \includegraphics[height=1cm]{./RDF_Logo.eps} & & \SetCell[c=3]{c} \large{Red de fríos S.A. de C.V.} & & & {\textbf{Código:}\\ \Codigo} \\
            \hline
            \SetCell[c=4]{c} {\Titulo} & & & & {\textbf{Fecha de publicación:}\\ \FechaPub} & {\textbf{Edición:}\\ \Edit} \\
            \hline
            \end{tblr}
	    }%
        \fancyfoot[CE,CO,LE,LO,RE,RO]{} %% clear out all footers
        \fancyfoot[C]{Página \thepage\ de \pageref{LastPage}}
        \fancyfoot[L]{\textbf{CONFIDENCIAL}}
        \fancyfoot[R]{\Codigo~v.\Edit}
}%

\setcounter{tocdepth}{4}
\setcounter{secnumdepth}{4}

\newcommand{\defglo}[1]{\item[\glsname{#1}] \glsdesc{#1}} % Define un comando para dar un item en un entorno de descripción para una palabra definida en el glosario


\begin{document}
\pagestyle{formato}
\include{src/PSA-ToC.tex} 
\thispagestyle{formato-PI}
\renewcommand{\MayorVer}{2}
\renewcommand{\MenorVer}{0}
\renewcommand{\Codigo}{PSA-1-PROG}
\renewcommand{\FechaPub}{2023--01}
%\renewcommand{\Edit}{2.1}
\renewcommand{\Titulo}{Inspección en la recepción de productos}

\section{\Titulo}\index{Inspección!en la recepción de productos}

\subsection{Objetivo}
\begin{itemize}
	\item \textbf{Establecer} un programa satisfactorio que asegure que la unidad de transporte se encuentra en condiciones adecuadas de mantenimiento, limpieza y temperatura apropiadas para el almacenamiento dentro de RDF.
\end{itemize}

\subsection{Alcance}
\begin{itemize}
	\item Este procedimiento aplica para todas las unidades con producto alimenticio recibidas en RDF;
	\item se extiende este documento al área de embarques, pero no se limita a otras áreas operativas.
\end{itemize}

\subsection{Términos y definiciones}


\begin{description}
	\defglo{producto}
	\defglo{producto-terminado}
	\defglo{alimento}
	\defglo{peligro-relacionado-con-la-inocuidad-de-los-alimentos}
\end{description}

\subsection{Documentos y/o normas relacionadas}
\begin{itemize}
	\item Programa de \glsfirst{BPD}.
\end{itemize}

\subsection{Procedimiento}\index{Procedimiento!inspección en la recepción de productos}

\subsubsection{Materiales}
\begin{itemize}
	\item Termómetro IR.
	\item Linterna.
\end{itemize}

\subsubsection{Precauciones de seguridad}

\begin{itemize}
	\item Usar \emph{uniforme completo} %TODO: Definición de uniforme completo.
\end{itemize}

\subsubsection{Instrucciones}

\paragraph{Inspección física del producto}\index{Instrucción!inspección fisica de productos}
\begin{enumerate}
	\item A la llegada del producto, se verifica la documentación del mismo. Si se descubre que el documento no cumple con los requisitos a la llegada, no se puede descargar el producto y el Supervisor de Almacén debe ser contactado inmediatamente para obtener instrucciones adicionales:
	      \begin{enumerate}
		      \item Documentos TIF \emph{(En caso de ser productos cárnicos);}
		      \item Listado de productos con cantidades;
		      \item Condiciones de almacenamiento.
	      \end{enumerate}
	\item Verificar la temperatura programada de la unidad con el producto entrante;
	\item Se revisa el número de sello que deberá coincidir contra la guía para asegurarse que cuadran. Si se descubre que el sello está roto a la llegada, no se puede descargar el producto y el Supervisor de Almacén debe ser contactado inmediatamente para obtener instrucciones adicionales;
	\item Se enrama la unidad para descargar el producto;
	\item La condición física del producto (caja, cubetas, etc.) debe ser revisado, cualquier daño debe ser documentado en el registro de entrada;
	\item Se verifica unidad interna. Cualquier evidencia de daño, suciedad, o plagas deben ser inmediatamente documento e informado al Supervisor de Almacén para obtener instrucciones adicionales antes de la descarga:
	      \begin{itemize}
		      \item Limpieza (suciedad, escombros, basura, etc.);
		      \item Evidencia de cualquier actividad de insectos (los excrementos, gomas de mascar, etc.);
		      \item Daños físicos (bolsas rotas, cajas dañadas, fugas, etc.);
		      \item Daño físico de la unidad (transporte);
		      \item Evidencia de violación de los materiales de empaque (caja o producto abierto);
		      \item presencia visible de agentes químicos.
	      \end{itemize}
	\item Verificar la temperatura del producto.\footnote{Todo producto refrigerado o congelado recibido en la instalación debe ser revisado para asegurar que fue enviado de origen a la temperatura adecuada, verificando la documentación correspondiente.}
	      \begin{itemize}
		      \item \textbf{Rangos de temperatura:}
		            \begin{itemize}
			            \item Los productos refrigerados DEBEN recibirse en un rango de temperatura de \qtyrange{2.7}{4.4}{\degreeCelsius}
			            \item Los productos congelados DEBEN recibirse a una temperatura \qty{<=-18}{\degreeCelsius}
		            \end{itemize}
		      \item De \textbf{no cumplirse} lo anterior se suspende el recibo por parte de personal de embarques y deberá notificarse inmediatamente al Supervisor del Almacén quien reportara al Gerente de Operaciones, quien a su vez informara de la situación al cliente y solicitara la ruta de acción.
	      \end{itemize}
	      % \begin{itemize}
	      % 	\item Se considerarán las especificaciones del producto para determinar si se aceptan o se rechazan los cargamentos.
	      % 	\item se comprobará que se cumple con la \gls{conformidad} de la temperatura 
	      % 	\item si es conforme la temperatura, se procederá a almacenar el producto en la camara designada según sus especificaciones.
	      % 	\begin{itemize}
	      % 	\item si el \gls{alimento} no requiere de condiciones de almacenamiento especiales\footnote{diferentes a las especificaciones genéricas de almacenamiento (ver \cref{esp:generica}).} se almacenará en la cámara adecuada. 
	      % 	\item si el \gls{alimento} requiere de condiciones de almacenamiento especiales y el cliente alquiló una camara exclusiva para su producto, entonces se almacenará ahí;
	      % 	\item de no ser así, se almacenará el \gls{alimento} en aquella camara compartida que coincida con el rango de las especificaciones del producto.
	      % 	\end{itemize}
	    %   \end{itemize}
	\item Todos los documentos de envío \emph{(números de lote, fechas, número de orden, etc.)} deberán ser registrados en el formulario de verificación de \gls{registro} y en la orden de entrada;
	\item Se almacena el \gls{alimento}.
\end{enumerate}

\paragraph{Identificación del producto}

\begin{enumerate}
	\item Al momento del envío de productos por parte de los proveedores se genera el documento de STOCK TRANFER el cual de manera inmediata es enviado a \gls{RDF} (control de Inventarios).
	\item Mesa de control recibe el documento \textit{STOCK TRANSFER} y genera el \emph{identificador de control de interno\footnote{Este es nombre que recibe el identificador de tarima.}} antes de la llegada del producto.
	\item La llegada del embarque a \gls{RDF} se tiene programada en fecha.
	\item Al momento de la llegada del embarque a \gls{RDF}, mesa de control coloca de manera inmediata a la descarga del producto la etiqueta identificador de control interno, al mismo tiempo que se realizará la verificación de cantidades de producto recibida
	\item El montacarguista debe trasladar el producto directamente del transporte a la cámara
\end{enumerate}

\paragraph{Reglamento de permanencia de producto en andén de carga}

\begin{itemize}
	\item Para \emph{\textbf{producto congelado:}} No más de \qty{1}{\hour};
	\item Para \emph{\textbf{productos refrigerados:}} No más de \qty{1}{\hour}.
\end{itemize}


\subsubsection{Responsables de la actividad}

\begin{itemize}
	\item \textbf{Ejecutado} por el personal de operaciones;
	\item \textbf{Monitoreado} por personal de aseguramiento de calidad;
	\item \textbf{Verificado} por personal de gerencia.
\end{itemize}

\subsubsection{Acciones preventivas}

\begin{itemize}
	\item Se llevará a cabo una inspección, registrando el indicador y medida correspondiente;
	\item si la desviación se repite, se capacitará al personal involucrado;
	\item Si después de la capacitación se repite una desviación, se tomarán acciones correctivas (\cref{sec:2.1:acc}).
\end{itemize}

\subsubsection{Acciones correctivas}
\label{sec:2.1:acc}
\begin{itemize}
	\item En caso contrario, la tarea se deberá volver a realizar como indica el procedimiento;
	\item En caso de no-conformidad, reportar en el formato de acciones correctivas.
\end{itemize}

\subsubsection{Frecuencia}

Cada recepción de \gls{alimento}

\begin{changelog}[simple, sectioncmd=\subsection*,label=changelog-1.2]
	\begin{version}[v=2.1, date=2023--01, author=Pablo E. Alanis]
		%\fixed
		\item Cambio de formato;
		\item Cambios en la serialización de versiones;
		\item Cambio de identificado, de PRO-OP-001 a OP-BPD-ESP-1.
		%\added
		\item Separación entre PPR y PPRO.
	\end{version}
\end{changelog}
\thispagestyle{formato-PI}
\renewcommand{\MayorVer}{2}
\renewcommand{\MenorVer}{0}
\renewcommand{\Codigo}{PSA-2-PRO} 
\renewcommand{\FechaPub}{2023--01}
%\renewcommand{\Edit}{2.1}
\renewcommand{\Titulo}{Procedimiento de lotificación de productos para su almacenamiento}

\section{\Titulo}\index{Procedimiento!lotificación de productos}

% \section{Procedimiento de lotificación de productos para su almacenamiento}

\subsection{Objetivo}

\begin{itemize}
	\item \textbf{Asegurar} que las tarimas con producto a almacenar sean identificadas con la información requerida por el cliente.
\end{itemize}

\subsection{Alcance}

\begin{itemize}
	\item Este documento es aplicable a todas las tarimas con producto que almacenadas en RDF.
\end{itemize}

\subsection{Términos y definiciones}
%TODO definiciones

\subsection{Procedimiento}

\subsubsection{Fundamento}

Los establecimientos y equipos dedicados al proceso de alimentos para consumo humano deben mantener los registros y el control de los productos y materiales de empaque por medio de la notificación de los mismos, con el cual es posible la realización de un sistema de rastreo desde su producción hasta su distribución.

\subsection{Materiales}

\begin{itemize}
	\item Factura de entrada
\end{itemize}

\subsection{Instrucciones}

\begin{itemize}
	\item Cada vez que se reciba en el área de almacén, productos o materiales de empaque, se debe registrar el lote que el cliente designado a este producto.
	\item El personal de recepción debe de verificar que el lote de los registros coincidan con el lote marcado en las tarimas del producto;
	\item Estos se darán de alta en el sistema de la empresa con el lote asignado por el proveedor o cliente.
	\begin{itemize}
		\item Si el producto no tiene lote o este no coincide se dará aviso al Supervisor de almacén (el lote asignado será asignado con el número de entrada)
	\end{itemize}
	\item Las tarimas que contienen el producto son marcadas con \emph{etiqueta de recepción de producto,} en ella se coloca la fecha de recepción, nombre del cliente, nombre del producto, fecha de caducidad, orden de compra o lotificación, numero de entrada, cantidad por tarima.
\end{itemize}

\subsection{Responsables de la actividad}

\begin{itemize}
	\item \textbf{Ejecutado} por personal de operaciones;
	\item \textbf{Monitoreado} por personal de calidad;
	\item \textbf{Verificado} por personal de gerencia.
\end{itemize}

\subsection{Acciones preventivas}

\begin{itemize}
	\item Se llevara a cabo una inspección. Registrando el indicador y medida correspondiente;
	\item Si la desviación se repite frecuentemente se dará curso de capacitación al personal de almacén y limpieza para que realicen eficientemente su trabajo;
	\item Si después de haber capacitado al personal de almacén se siguen presentando desviaciones por causas injustificadas, se tomaran acciones más enérgicas con el personal por incumplimiento con sus deberes.
\end{itemize}

\subsection{Acciones correctivas}

\begin{itemize}
	\item En caso contrario la tarea se deberá volver a realizar como se indica en el procedimiento.
	\item En caso de no conformidad reportar en \RAC.
\end{itemize}

\subsection{Frecuencia}

\begin{itemize}
	\item Cada recepción de Producto
	\item Cada entrega de producto.
\end{itemize}

\begin{changelog}[simple, sectioncmd=\subsection*,label=changelog-2.2]
	
	\begin{version}[v=2.1, date=2023--01, author=Pablo E. Alanis]
			\item Cambio de formato;
			\item Cambios en la serialización de versiones;
	\end{version}

	\begin{version}[v=1.3, date=2022--05, author=Alonso M.]
		\item cambio de fecha;
		\item cambio de código.
	\end{version}

	\shortversion{v=1.2, date=2021--05, changes=No hubo cambios}
\end{changelog}
\renewcommand{\MayorVer}{2}
\renewcommand{\MenorVer}{1}
\renewcommand{\Codigo}{PSA-1-PROG}
\renewcommand{\FechaPub}{2023--01}
%\renewcommand{\Edit}{2.1}
\renewcommand{\Titulo}{Inspección de producto durante su almacenamiento}

\section{\Titulo}
\index{Inspección!de productos durante su almacenamiento}

\subsection{Objetivo}

Establecer un programa de inspección que asegure que las condiciones de almacenamiento se encuentren en condiciones óptimas de acuerdo a un buen mantenimiento de instalaciones, limpieza y temperatura apropiada para la conservación de productos alimenticios.

\subsection{Alcance}

A todos las áreas de almacenamiento de producto

\subsection{Terminología y definiciones}

N/A

\subsection{Documentos y/o normas relacionados}

\begin{itemize}
	\item 1.0 - BPD - Indice|Programa de buenas prácticas de manufactura
\end{itemize}

\subsection{Procedimiento}

\subsubsection{Materiales}

\begin{itemize}
	\item N/A
\end{itemize}

\subsubsection{Precauciones de Seguridad}

\begin{itemize}
	\item Usar uniforme completo.
\end{itemize}

\subsection{Instrucciones}

\subsubsection{Temperatura de Almacenaje de productos Refrigerados y Congelados.}

\begin{enumerate}
	\item Para asegurar que todo producto refrigerado o congelado se mantenga a la temperatura adecuada durante el almacenamiento, todos los productos se deben colocar directamente después de ser recibido en la cámara de refrigeración o congelación según la asignación, una vez que se descarguen.
	\item Las cámaras de refrigeración y de congelación están equipadas con un termómetro calibrado (RACK), en el cual se puede verificar la temperatura de cada cámara.
	\item Las lecturas de temperaturas deben ser documentadas tres veces durante el turno los días laborados. El Supervisor de mantenimiento es responsable de la verificación de temperaturas.
	\item Las puertas del andén de embarques deben permanecer cerradas en todo momento a menos que la unidad este siendo utilizada para carga y descarga de productos.
	\item Las puertas del andén deben de ser abiertas cuando la unidad de transporte ya se encuentre colocada en la rampa, esto con el fin de evitar fuga de temperatura y / o entrada de polvo o plaga.
	\item Cualquier daño a las cámaras de refrigeración o congelación (techo, paredes, puertas, cortinas, etc.) debe ser reparadas en el menor tiempo posible. El Supervisor de mantenimiento es responsable de garantizar que todo el equipo este en buen estado y funcionando adecuadamente.
\end{enumerate}

\subsubsection{Almacenaje y manejo de Productos Secos / Ambiente, Refrigerados y Congelados}

\begin{enumerate}
	\item Para asegurarse de que los productos ambiente, refrigerados, y congelados son almacenados y protegidos de daños físicos, químicos, biológicos (contaminación), todos los productos deben ser almacenados 30 centímetros desde todas las paredes interiores (si el diseño del almacén lo permite), sobre atados limpios, 30 centímetros de las superficies del techo, así como lejos de todas las luces superiores y nunca almacenados directamente en contacto con los pisos (debe de estar sobre tarima)
	\item Todos los productos deben guardarse en su empaque secundario y etiquetados con las fecha de recepción necesarias.
	\item Todos los productos deben ser rotados de acuerdo en el principio de “First-in-First-Out” \textbf{(FIFO)} o primeras entradas primero a vencer (O según instrucciones de cada cliente).
	\item Un área de “retención” debe haber en el área de almacenamiento para almacenar cualquier artículo recuperado, muestras de productos, producto dañado.
	\item Todos los productos deben mantenerse libres de polvo y suciedad en todo momento.
	\item Cualquier o todos los productos químicos almacenados en la instalación deben estar completamente separados de todos los envases y productos alimenticios – SIN EXCEPCIONES. Todos los productos químicos ubicado en el centro debe tener una hoja de identificación (NO es el caso para este almacén).
\end{enumerate}

\subsubsection{Control y mantenimiento de temperatura}

\begin{enumerate}
	\item Es imperativo que las cámaras en refrigerado y congelado en nuestras instalaciones se mantengan a la temperatura adecuada en todo momento. Esto incluye horas regulares, después de horas regulares, fines de semana y días festivos.
	\item Existe un procedimiento para el seguimiento de temperaturas del refrigerador y del congelador durante el horario normal \%\%CUAL?\%\%las temperaturas son supervisadas y documentadas al día en tres tiempos (Inicio de turno, medio turno y fin de turno) de Lunes a Sábado. El Supervisor de Mantenimiento es inmediatamente informado de cualquier irregularidad detectada en la temperatura (incluyendo horas después del trabajo, fines de semana y días festivos), de modo que las acciones correctivas se puedan tomar.
\end{enumerate}


\subsubsection{Plan alterno en caso de avería de equipos de enfriamiento con productos}

Al fin de garantizar que todos los productos refrigerados y congelados estén protegidos contra cualquier peligro microbiológico asociado con la perdida de temperatura, las cámaras con refrigeración o congelación que presenten una falla y que no puedan ser reparadas en menos en el caso de refrigeración de 4 horas y en congelación de 10 hrs (siempre y cuando las puertas de cámara se encuentren cerradas y el producto mantenga la temperatura indicada), el producto DEBE ser trasladado inmediatamente a otra cámara (previamente enfriada).

\subsection{Responsables de la actividad}

\begin{itemize}
	\item \textbf{Ejecutado} por personal de operaciones y mantenimiento;
	\item \textbf{Monitoreado} por personal de calidad;
	\item \textbf{Verificado} por personal de gerencia.
\end{itemize}

\subsection{Acciones preventivas}

\begin{itemize}
	\item Se llevara a cabo una inspección. Registrando el indicador y medida correspondiente
	\item Si la desviación se repite frecuentemente se dará curso de capacitación al personal de almacén y limpieza para que realicen eficientemente su trabajo.
	\item Si después de haber capacitado al personal de almacén se siguen presentando desviaciones por causas injustificadas, se tomaran acciones más enérgicas con el personal por incumplimiento con sus deberes.
\end{itemize}

\subsection{Acciones correctivas}

\begin{itemize}
	\item En caso contrario la tarea se deberá volver a realizar como se indica en el procedimiento.
	\item En caso de no conformidad reportar en formato de acciones correctivas F-OP-40.xls
\end{itemize}

\subsection{Frecuencia}

\begin{itemize}
	\item \textbf{Verificación de temperatura:} tres veces al día.
\end{itemize}

\subsection{Historial de modificaciones}

\begin{itemize}
	\item \textbf{Cuarta edición:} cambio de fecha de 03 de noviembre 2017 a 28 de enero 2019, se realizó cambio de código de OA-P-IDA a PR-008 y no se realizaron cambios de la revisión 003 a la 004.
	\item \textbf{Quinta edición:} febrero 2020 se hizo cambio de formato y cambio de código de PR-008 a PRO-OP-005. Y no se realizaron cambios de la revisión 04 a la 05.
	\item \textbf{Sexta edición:} febrero 2021 no se realizaron cambios de la revisión 05 a la 06.
	\item \textbf{Séptima edición:} febrero 2022 no se realizaron cambios de la revisión 06 a la 07.
	\item \textbf{Octava edición:} 2023--01: cambios en serialización y en formato.
\end{itemize}

\renewcommand{\MayorVer}{2}
\renewcommand{\MenorVer}{1}
\renewcommand{\Codigo}{PSA-1-PROG}
\renewcommand{\FechaPub}{2023--01}
%\renewcommand{\Edit}{2.1}
\renewcommand{\Titulo}{PEPS o VENCER}

\section{\Titulo}
\index{Procedimiento!PEPS o VENCER}

% \section{PEPS o VENCER}

\subsection{Objetivo}
Asegurar que los productos, serán en sus totalidades lotificadas e identificados y verificados sus movimientos por medio de conteos programados de acuerdo a lo establecido en este procedimiento.

\subsection{Alcance}
Todo producto que ingrese a las instalaciones del almacén.

\subsection{Terminología y definiciones}
Inventario de suministros y materiales. La percepción de inventario asegura que la compañía tiene materiales a la mano para hacer productos y que los fondos no son desperdiciados en materiales innecesarios. Una cuenta precisa de inventario también permite que las compañías controlen y ordenen suficientes, materiales para la demanda de los consumidores.

\subsection{Documentos y/o normas relacionadas}
\begin{itemize}
	\item Registros de recepción
\end{itemize}

\subsection{Procedimiento}

\subsubsection{Materiales}

\begin{itemize}
	\item N/A
\end{itemize}

\subsubsection{Precauciones de seguridad}

\begin{itemize}
	\item Usar el uniforme completo.
\end{itemize}

\subsubsection{Instrucciones}

\paragraph{Primeras entradas, primeras salidas (PEPS) o VENCER}

\begin{enumerate}
	\item Hacer la revisión correspondiente en cada recepción de producto, la revisión abarca todos los identificadores únicos del envió (número de lote, fechas, cantidades) deben estar documentados.
\end{enumerate}


\begin{enumerate}
	\item Se etiqueta por tarima y por lote cada producto;
	\item La etiqueta de \emph{"Información de material"} contiene los siguientes datos:
	\begin{itemize}
		\item Nombre de producto
		\item Nombre del cliente
		\item Fecha de recepción
		\item Identificador de entrada
		\item Lote
		\item Cantidad
	\end{itemize}
\end{enumerate}

\subparagraph{En operación}

\begin{itemize}
	\item Cada insumo con fecha de caducidad más larga (Producto nuevo), se acomodara en la parte superior trasera de cada rack y los insumos que ya estaban presentes se recorren a la parte inferior frontal.
	\item Al surtir los insumos, el personal de almacén deberá entregar los que tengan la fecha de recepción más antigua o el producto que le cliente asigne para su distribución.
\end{itemize}

\subparagraph{En sistema}

\begin{itemize}
	\item En el sistema de inventarios se cargan las entradas de producto en general al momento de su llegada, en este sistema administra el orden de salida de los mismos y permite dar de baja únicamente en el orden de primeras entradas primeras salidas.
	\item El cliente también pudiera dar instrucciones de cual producto dar salida de forma primaria por estrategia comercial.
	\item En ambos casos el sistema de manera automática da de baja del sistema el producto que físicamente se da salida, actualizando las existencias.
\end{itemize}

\subsection{Responsables de la actividad}

\begin{itemize}
	\item \textbf{Ejecutado} por personal de operaciones;
	\item \textbf{Monitoreado} por personal de calidad;
	\item \textbf{Verificado} por personal de gerencia.
\end{itemize}

\subsection{Acciones preventivas}

\begin{itemize}
	\item Se llevara a cabo una inspección. Registrando el indicador y medida correspondiente
	\item Si la desviación se repite frecuentemente se dará curso de capacitación al personal de almacén y limpieza para que realicen eficientemente su trabajo.
	\item Si después de haber capacitado al personal de almacén se siguen presentando desviaciones por causas injustificadas, se tomaran acciones más enérgicas con el personal por incumplimiento con sus deberes.
\end{itemize}

\subsection{Acciones correctivas}

\begin{itemize}
	\item En caso contrario la tarea se deberá volver a realizar como se indica en el procedimiento.
	\item En caso de no conformidad reportar en formato de acciones correctivas F-OP-40.xls
\end{itemize}

\subsection{Frecuencia}

\begin{itemize}
	\item Diaria.
\end{itemize}

\subsection{Historial de modificaciones}

\begin{itemize}
	\item \textbf{Cuarta edición:} cambio de fecha de 03 de noviembre 2017 a 28 de enero 2019, se realizó cambio en el código de CA-P-PEPV a PR-009. Y no se realizaron modificaciones de la revisión 03 a la 04
	\item \textbf{Quinta edición:} febrero 2020 se hizo cambio de forma y cambio de código de PR-009 a PRO-OP-006. Y no se realizaron modificaciones de la revisión 04 a la 05.
	\item \textbf{Sexta edición:} febrero 2021 no se realizaron modificaciones de la revisión 05 a la 06.
	\item \textbf{Séptima edición:} febrero 2022 no se realizaron cambios de la revisión 06 a la 07.
\end{itemize}

\subsection{Listado de distribución}

N/A

\subsection{Anexos}

N/A

%\thispagestyle{formato-PI}
\renewcommand{\TipoID}{CI}
\section{Evaluación de proveedores}
\renewcommand{\Codigo}{\Prog--\thesection--\TipoID}
Estimado proveedor de Red de Fríos, S.A. de C.V., el presente formato de evaluación tiene por objetivo conocer los procesos involucrados en el surtimiento del insumo que nos provee y forma parte de nuestro Sistema de Calidad de los Alimentos. Con la información que nos proporcione realizaremos una evaluación y dependiendo de los resultados le calificaremos como un \emph{proveedor aprobado, aprobado condicionado o rechazado,} de los resultados de esta evaluación probablemente se desprendan acciones que le solicitaremos sean implementadas para mejorar la calidad del servicio que nos brinda. Agradecemos de antemano su apertura y las facilidades que nos brinda para realizar esta evaluación.

\subsection{Datos generales del proveedor}


\subsection{Cumplimiento de especificaciones}

Se han establecido con el proveedor las especificaciones del insumo requerido y el proveedor cumple con estas especificaciones en el 100\% de las entregas a RDF.

\subsection{Cumplimiento de fechas y cantidades comprometidas}

El proveedor \emph{cumple con las fechas comprometidas de entrega} de insumos así como con las cantidades requeridas en la orden de compra.

\subsection{Capacidad para surtir cantidades requeridas}

El proveedor cuenta con la \emph{capacidad para surtir las necesidades} de volumen de producto que RDF le solicita.

\subsection{Capacidad para surtir cantidades extraordinarias}

El proveedor tiene la capacidad de responder en tiempo razonable a demandas extraordinarias de volumen del producto que nos provee.

\subsection{Condiciones de producción y almacenamiento}

El proveedor mantiene un ambiente de producción y/o almacenamiento tal que no se observan condiciones que representen un riesgo físico, químico o biológico a la inocuidad de los alimentos almacenados en nuestras instalaciones.

\subsection{Presencia de plagas}

Durante la visita de las instalaciones no se observa infestación por algún tipo de plaga y preferentemente el proveedor cuenta con un sistema de control de plagas.

\subsection{Mantenimiento de equipos e instalaciones}

Durante la visita de las instalaciones se observan las instalaciones y maquinaria bien mantenidas de tal forma que se asegura su buen funcionamiento y por lo tanto el abasto del insumo.

Agradecemos nuevamente las facilidades brindadas para realizar esta evaluación, en los próximos días recibirá el resultado final. En este mismo reporte recibirá toda aquella solicitud de acción a implementar en sus instalaciones que haya sido mencionada durante el recorrido; le pedimos responder por escrito a estas solicitudes mencionando el responsable de implementarlas y la fecha compromiso para su finalización.


Nombre y firma del proveedor


Nombre y firma de red de fríos

\thispagestyle{formato-PI}
\renewcommand{\MayorVer}{2}
\renewcommand{\MenorVer}{0}
\renewcommand{\Codigo}{PSA-6-CE}
\renewcommand{\FechaPub}{2023--01}
\renewcommand{\TipoID}{PRO}
\renewcommand{\Titulo}{Procedimiento para solicitud de citas}

\section{\Titulo}\index{Procedimiento!solicitud de citas}
\renewcommand{\Codigo}{\Prog--\thesection--\TipoID}
Estimado cliente, a continuación le presentamos el proceso y lineamientos para la solicitud de citas tanto ingresos como salidas de material.

Le solicitamos que las citas sean solicitadas mínimo con 24 horas de anticipación mediante correo electrónico al equipo de mesa de control en los siguientes horarios:

\begin{itemize}
	\item \textbf{Lunes a viernes de 08:00 am a 16:00 pm}
	\item \textbf{Sábados de 08:00 am a 11:30 am.}
\end{itemize}

Quedando las solicitudes de la siguiente forma:

\definecolor{Gallery}{rgb}{0.929,0.929,0.929}
\begin{longtblr}[
	label = citas:solucitud,
	entry = Procedimiento de solicitud de citas.,
	caption = Horarios para solicitud de citas.
	]{%
	width = \linewidth,
	colspec = {Q[125]Q[129]Q[108]Q[577]},
	cells = {c},
	row{even} = {Gallery}
	}
	\toprule
	Día de solicitud & Hora límite para solicitud & Día de carga & Observación                                                                         \\
	\midrule
	Lunes            & 15:30 pm                   & martes       & La cita debe realizarse antes de las 15:30 pm para que se asigne cita el Martes.    \\
	Martes           & 15:30 pm                   & miércoles    & La cita debe realizarse antes de las 15:30 pm para que se asigne cita el miércoles. \\
	Miércoles        & 15:30 pm                   & jueves       & La cita debe realizarse antes de las 15:30 pm para que se asigne cita el jueves.    \\
	Jueves           & 15:30 pm                   & viernes      & La cita debe realizarse antes de las 15:30 pm para que se asigne cita el viernes.   \\
	Viernes          & 15:30 pm                   & sábado       & La cita debe realizarse antes de las 15:30 pm para que se asigne cita el sábado.    \\
	Sábado           & 11:30 am                   & lunes        & La cita debe realizarse antes de las 11:30 am para que se asigne cita el lunes.     \\
	\bottomrule
\end{longtblr}

\begin{itemize}
	\item \textbf{Los pedidos enviados fuera de estos horarios quedarán a disposición de andén;\footnote{En espera de que tengamos una ventana disponible.}}
	\item Es \emph{requisito indispensable} presentar el correo de confirmación impreso, que especifica el número de pedido a surtir y la hora asignada a su cita;
	\item Es importante que en la solicitud de la cita nos indiquen el \emph{nombre completo de la persona que realizará la recolección o entrega,} así como los datos de la unidad:
	      \begin{itemize}
		      \item Línea de transporte;
		      \item placas;
		      \item número de tractor y placas;
		      \item número de contenedor.
	      \end{itemize}
	\item Nuestro horario de atención sin cargo de tiempo extra es de 8:00 a 17:00 de lunes a viernes. El camión a carga o descarga debe presentarse una hora antes de las 17:00 para que no se extienda su servicio; lo mismo aplica para los sábados de 8:00 a 13:00; la última cita se otorga a las 16:00 de lunes a viernes y los sábados a las 12:00.
\end{itemize}

\newpage

\subsection{Contactos Red de Fríos}

\begin{contact}[Mesa de control] \label{contact:MesaDeControl}
	Magdalena Aguilar Robledo\\
	Teléfono: \textbf{(81) 8351-6685 ext. 105}\\
	\email{mesacontrolrayon@reddefrios.com}
\end{contact}

\begin{contact}[Jefe de almacén] \label{contact:JefeDeAlmacen}
	Luis Ángel Escareño Aguirre\\
	Teléfono: \textbf{(81) 8351-6685 ext. 105}\\
	Celular: \textbf{811-390-7099}\\
	\email{lescareno@reddefrios.com}
\end{contact}

\begin{contact}[Gerente de operaciones] \label{contact:Gerente}
	Gerardo Ruiz Torres, QFB\\
	Teléfono: \textbf{(81) 8351-6685 ext. 102}\\
	Celular: \textbf{818-029-5609}\\
	\email{gruiz@reddefrios.com}
\end{contact}

\begin{contact}[Médico TIF] \label{contact:MedicoTIF}
	M.V.Z. Rocío Luna\\
	Teléfono: \textbf{(81) 8351-6685 ext. 103}\\
	\email{luna\_mvz\_ceu@live.com.mx}
\end{contact}

Es necesario indicar:
\begin{enumerate}
	\item Nombre del producto a guardar o retirar;
	\item factura que lo ampara;
	\item código con el que usted lo maneja\footnote{Se dará de alta el mismo en nuestro sistema};
	\item lote;
	\item fecha de caducidad;
	\item cantidad de tarimas y/o cajas que se guardarán o retirarán;
	\item especificaciones de temperatura de almacenamiento.
	\item alérgenos presentes en el \gls{alimento} (ver \cref{nota:Alergenos})
\end{enumerate}

\begin{note}[Listado de alérgenos según la NOM-051-SCFI/SSA1-2010] \label{nota:Alergenos}
	Es de suma importancia que nos indique si el producto a resguardar es un Alérgeno o contiene algún ingrediente Alérgeno en su composición.
	\begin{itemize}
		\item Leche
		\item Trigo
		\item Frutas Secas
		\item Huevos
		\item Soya y derivados
		\item Cacahuates
		\item Pescados
		\item Crustáceos
		\item sulfitos en concentración \qty{\geq 10}{\milli\gram\per\kilo\gram}.
	\end{itemize}
\end{note}

Para carga o descarga requerimos contar con los datos de la unidad, placas, nombre del chofer y línea de transporte.

Para ingresos de mercancía es necesario presentar la siguiente documentación:

\begin{itemize}
	\item Factura, remisión y pedimento (para importados) del producto a recibir.
	\item Certificado de calidad de los productos a recibir.
	\item Certificado de fumigación de la unidad de transporte.
\end{itemize}

En caso de ser \emph{producto cárnico} es necesario presentar \emph{Aviso de Movilización TIF} con destino a nuestro almacén \emph{TIF 278.}

Para cuestión de \emph{Pescados y Mariscos} es necesario presentar la \emph{Guía de Pesca} del producto con destino a nuestro almacén TIF 278.

Para la \emph{salida de producto TIF} se requiere se nos informen los datos de la unidad que lo cargará:
\begin{itemize}
	\item Línea de Transporte;
	\item chofer;
	\item placas;
	\item marca del contenedor o caja;
	\item número de destino TIF.
\end{itemize}

\begin{note}[Datos de aviso TIF]\label{note:DatosAvisoTIF}
	TIF 278\\
	Red de Fríos S.A. de C.V.\\
	Av. I. López Rayón \# 2810.\\
	Col. Bella Vista.\\
	C.P. 64410, Monterrey N.L., México
\end{note}

\begin{changelog}[title=Registro de cambios,simple, sectioncmd=\subsection*,label=changelog-\thesection-\MayorVer.\MenorVer]
	\begin{version}[v=\MayorVer.\MenorVer, date=2023--01, author=Pablo E. Alanis]
		\item Cambio de formato;
		\item los contactos se agregaron a recuadros;
		\item cambio de código;
		\item se agregó recuadro sobre alérgenos.
	\end{version}

	\begin{version}[v=1.7, date=2022--05, author=Alonso M.]
		\item No se realizaron modificaciones de la revisión 05 a la 06.
	\end{version}
\end{changelog}
\thispagestyle{formato-PI}
\renewcommand{\MayorVer}{1}
\renewcommand{\MenorVer}{0}
\renewcommand{\FechaPub}{2023--01}
\renewcommand{\TipoID}{G}
\renewcommand{\Titulo}{Programa de Seguridad Alimenticia --- Glosario}
\renewcommand{\Codigo}{\Prog--\TipoID}
% \printunsrtglossary[type=acronym,nonumberlist,style=long]
% \printunsrtglossary[type=\glsdefaulttype]
\printglossary
% \printunsrtglossary[type=acronym,nonumberlist]
% \printunsrtglossary[type=\glsdefaulttype,nonumberlist]
\thispagestyle{formato-PI}
\renewcommand{\MayorVer}{1}
\renewcommand{\MenorVer}{0}
\renewcommand{\FechaPub}{2023--01}
%\renewcommand{\Edit}{01}
\renewcommand{\TipoID}{I}
\renewcommand{\Titulo}{Programa de Buenas Prácticas de Distribución --- Indice alfabético}
\renewcommand{\Codigo}{\Prog--\TipoID}
\printindex
\end{document}